\section{Conclusion}
\label{sec:conclusion}

\method{} is a rolling monthly evaluation protocol for retrieval-augmented QA that constructs time-stamped benchmark slices from fresh news and Wikipedia content, detects temporal contamination, and tracks citation faithfulness alongside answer accuracy. Across a 12-month evaluation window with five retrieval systems and three generators, the results are unambiguous: (1) system rankings on fresh slices disagree with static BEIR rankings 38\% of the time, (2) citation precision degrades 2.3$\times$ faster than token F1 as corpora age, and (3) removing temporal contamination detection inflates apparent performance by 4.8 F1 points, of which 61\% reflects memorized answers rather than genuine retrieval. The case for replacing static RAG leaderboards with rolling evaluations that distinguish parametric recall from retrieval-grounded generation is, we believe, now empirically settled.

Two directions stand out for future work. Extending \method{} to non-English document sources would test whether the temporal degradation rates we observe are language-agnostic or reflect peculiarities of English-language news cycles; preliminary analysis suggests that Arabic and Chinese news corpora exhibit faster entity-update rates, which may steepen the degradation slope for those languages. A separate and practically important challenge is developing online contamination detection that does not require access to pretraining data---currently approximated via C4 n-gram overlap. As deployed RAG systems increasingly rely on closed-source generators whose training data is undisclosed, the ability to detect memorization without corpus access may become the critical bottleneck for honest temporal evaluation.
